\documentclass[11pt,a4paper]{article}

% ---------------------------------------------------------------------------
% ZOE Platform North Corfu — Individual Report (Anshul Agrawal)
% DSR Phase 4: Demonstration. English, APA7. Compiles with pdflatex.
% Image placeholders are marked with: % TODO: Bild einfügen — <description>
% Citations are drawn EXCLUSIVELY from docs/literature-index.md.
% ---------------------------------------------------------------------------

\usepackage[T1]{fontenc}
\usepackage[utf8]{inputenc}
\usepackage[margin=1in]{geometry}
\usepackage{graphicx}
\usepackage{xcolor}
\usepackage{enumitem}
\usepackage{titlesec}
\usepackage[hidelinks]{hyperref}

\setlength{\parindent}{0pt}
\setlength{\parskip}{0.6em}
\titlespacing*{\section}{0pt}{1.2em}{0.5em}
\titlespacing*{\subsection}{0pt}{0.9em}{0.3em}

% APA7 hanging-indent reference entry
\newcommand{\refitem}[1]{\noindent\hangindent=2em\hangafter=1 #1\par\smallskip}

% Image / screenshot placeholder box
\newcommand{\figplaceholder}[1]{%
  \begin{center}
  \fbox{\parbox[c][2.6cm][c]{0.82\textwidth}{\centering\itshape [Placeholder --- #1]}}
  \end{center}}

\title{\textbf{Demonstrating a Civic Sustainability Platform for the\\ Municipality of North Corfu}\\[0.3em]
\large A Design Science Research Account of Phase~4 (Demonstration)}
\author{Anshul Agrawal\\ \small Project Seminar Information Systems, FAU Erlangen--N\"urnberg, Group~1 (SoSe~2026)}
\date{June 2026}

\begin{document}
\maketitle

\begin{abstract}
\noindent This report documents Phase~4 (\emph{Demonstration}) of a Design Science Research (DSR) project that develops ZOE, a digital sustainability and civic-participation platform for the Greek municipality of North Corfu. Building on the problem identification and solution objectives elaborated by the team (Phases~1--2) and the design principles developed by my project partner (Phase~3), this report describes the implemented artefact and demonstrates how it addresses the identified sub-problems through a set of scenario walkthroughs. Each scenario is linked to a design principle and to the literature that justifies it, classified as evidence of a measured effect (Type~A) or as a framework recommendation (Type~B). The report further argues for the \emph{evaluability} of the artefact, showing that artificial, ex-ante evaluation (WCAG conformance, Lighthouse performance) is already feasible, while naturalistic, ex-post evaluation is planned following the Framework for Evaluation in Design Science (FEDS; Venable et al., 2016). The contribution is positioned as an \emph{improvement / exaptation} (Gregor \& Hevner, 2013): the transfer of established e-participation and transparency designs into the specific, under-served context of a small Mediterranean island municipality.
\end{abstract}

\section{Introduction}

North Corfu is a small municipality at the northern tip of the Greek island of Corfu, in the Ionian Sea. Its economy and identity are shaped by a tension that is visible across the Mediterranean: an exceptional natural environment --- Natura~2000 wetlands such as the Antinioti Lagoon, the forested Erimitis peninsula, and a 22~km coastline --- set against an intense tourism pressure of roughly four million annual visitors, whose seasonal waste once overwhelmed the island's only landfill at Temploni in 2018. The municipal \emph{ZOE} programme bundles a portfolio of environmental actions (among them sea-turtle protection, beach cleanups carried out with schools, and the protection of natural monuments) and frames them against the seventeen United Nations Sustainable Development Goals (SDGs).

The practical trigger for this project is a communication deficit. At present, the municipality's environmental actions are communicated largely through the personal Facebook profile of the vice-mayor; there is no central, structured, multilingual place where residents, visitors, schools, and researchers can find the actions, understand their contribution to the SDGs, see the impact of past participation, or contribute their own initiatives. This is not a peculiarity of Corfu: empirical analyses of Greek municipal websites find that none scored above 55\% on a combined performance and user-experience audit (Tsatsani et al., 2024), and a survey of the Greek smart-city transition reports high public awareness (82.8\%) alongside considerable barriers and dissatisfaction with existing infrastructure (Carayannis et al., 2026).

The problem space --- the sub-problems and the solution objectives --- was elaborated collectively by the project team during Phases~1 and~2 and is therefore only summarised here; the reader is referred to the contributions of my team colleagues for its full derivation. In condensed form, the team identified the following sub-problems (TP): a visibility deficit (TP1), a participation barrier (TP2), a lack of SDG transparency (TP3), the heterogeneity of target groups (TP4), and tourists as an un-tapped resource rather than only a burden (TP6). The translation of these problems into design principles is the subject of my partner's report (Phase~3).

This report concentrates on \textbf{Phase~4, Demonstration}. Its research questions are: \emph{(RQ1)} How can the demonstration phase of a DSR project show that the ZOE artefact addresses the identified sub-problems in realistic usage scenarios? and \emph{(RQ2)} How does demonstration differ from, yet prepare, the evaluation of the artefact (Phase~5)? The objective is therefore twofold: to describe the implemented instantiation and to demonstrate it through scenario walkthroughs that each make a design principle observable in practice, while carefully separating what was \emph{built} from what was only \emph{conceived}.

\section{Research Method: Design Science Research}

\subsection{DSR in general}

Design Science Research is a paradigm in which knowledge is created through the construction and evaluation of \emph{artefacts} that are intended to solve a class of problems. The canonical process model of Peffers et al. (2007) structures a DSR project into six activities: (1)~problem identification and motivation, (2)~definition of the objectives of a solution, (3)~design and development, (4)~demonstration, (5)~evaluation, and (6)~communication. Crucially, the model is iterative and admits several entry points, so that the six activities are better understood as a cycle than as a strict waterfall.

Hevner (2007) complements this process view with a structural one, describing DSR as three interlocking cycles. The \emph{Relevance Cycle} connects the research to its application environment, drawing requirements from it and returning artefacts to it for field testing. The \emph{Rigor Cycle} connects the research to the knowledge base, grounding the work in existing theories and methods and feeding new knowledge back. The central \emph{Design Cycle} iterates tightly between construction and evaluation. For the present project, North Corfu supplies the relevance, the curated literature corpus supplies the rigor, and the documented iterations (recorded in the project's development log) constitute the design cycle.

The form that design knowledge takes is specified by Gregor and Jones (2007), who identify eight components of a design theory, including its \emph{principles of form and function}, its \emph{principles of implementation}, and an \emph{expository instantiation}. The last of these is directly relevant to a demonstration: an instantiation is the concrete software artefact through which an otherwise abstract design principle becomes observable. Finally, Gregor and Hevner (2013) provide a vocabulary for the \emph{knowledge contribution} of a DSR project, distinguishing, among others, \emph{improvement} (a new solution for a known problem) and \emph{exaptation} (the transfer of a known solution into a new domain).

\subsection{Positioning of Phase~4}

Within this framework, the present report addresses Phase~4. Demonstration, in the sense of Peffers et al. (2007), is the use of the artefact to solve one or more instances of the problem; it is the empirical-yet-not-yet-evaluative step in which the instantiation is exercised in representative situations. It follows the design and development of the artefact (Phase~3, the focus of my partner's report) and precedes its systematic evaluation (Phase~5). Methodologically, demonstration corresponds to the \emph{expository instantiation} of Gregor and Jones (2007): it shows that the design principles can be realised and that they behave as intended.

It is important to keep the proportion of this report appropriate. The general method is a foundation, not the main contribution; in line with the project's reporting convention, the methodological chapter is deliberately kept to roughly half the length of the main chapter (Sections~4 below).

\section{Theoretical Background and Related Work}

The knowledge base for the ZOE artefact spans three fields, each of which supplies designs that the demonstration exercises.

\textbf{Comparable municipal e-participation.} The most widely cited reference platform is Decide Madrid, built on the open-source Consul software, which has been adopted in more than one hundred institutions across thirty-three countries. Royo et al. (2020) analyse it as a critical case and find that, while high citizen expectations explain the initially high participation, a lack of transparency and poorly functioning participatory features degrade its performance over time. Arana-Catania et al. (2021) identify information overload as a central barrier on the same platform and demonstrate that natural-language-processing techniques can mitigate it. Royo et al. (2024) compare Decide Madrid with the Scottish ``We Asked, You Said, We Did'' initiative and emphasise the importance of \emph{closing the feedback loop}. These works are Type~A evidence: they report measured or case-derived effects, and they motivate ZOE's emphasis on curation and on making impact visible.

\textbf{Mediterranean and tourism analogues.} Vegas Macias et al. (2023) study the fishing village of Marsaxlokk in Malta and show, through action research, how a digital storytelling platform let a fishing community co-create heritage-tourism experiences and raised awareness for sustainable tourism --- a setting strikingly analogous to North Corfu. Laksmi et al. (2026) provide stronger, quasi-experimental evidence from sixty tourism destinations in Bali: a combination of community participation and digital strategy significantly improved destinations' sustainability performance, with community participation and digital strategy emerging as the two pivotal factors. Both are Type~A and directly justify ZOE's treatment of visitors as a resource (TP6).

\textbf{The Greek and accessibility context.} Tsatsani et al. (2024) and Carayannis et al. (2026) ground the work in the Greek reality of constrained municipal digital services and citizen attitudes. On accessibility, Csontos and Heckl (2021, 2025) show that public-sector websites routinely fail to meet the WCAG~2.1 AA level, and Pontus and Rodr\'iguez V\'azquez (2021) show that localised language versions are frequently \emph{less} accessible than the original-language version --- a finding that directly shapes how ZOE's demonstration must treat its three languages.

Methodologically, the demonstration is an expository instantiation (Gregor \& Jones, 2007), and the contribution is best read as improvement / exaptation (Gregor \& Hevner, 2013): designs proven elsewhere are adapted to an under-served small-municipality context.

\section{Implementation and Demonstration of the Artefact}

This is the central chapter. It first summarises what was implemented, separating built from conceived features, and then demonstrates the artefact through four scenario walkthroughs. Each scenario is described in text with a placeholder for the corresponding screenshot or live-demo step, and each is tied to a design principle (DP, defined in the project matrix) and to its justifying source.

\subsection{The implemented instantiation (built vs.\ conceived)}

ZOE is a full-stack web application. The frontend is built with React~19, TypeScript (strict), Vite, Tailwind CSS, and React Router; internationalisation uses \texttt{react-i18next} for English, Greek, and German; global state uses Zustand. A Node.js/Express/Prisma/SQLite backend with JWT authentication supports an administrative area for the municipality. Quality is safeguarded by a Vitest unit/integration suite and a Playwright end-to-end suite, including automated accessibility checks with axe-core.

\textbf{Built and tested.} A central public catalogue of actions with filtering and a per-action detail page (TP1); an SDG dashboard and per-action SDG badges, and a transparency page with key performance indicators and progress bars (TP3); a low-threshold participation page with five contribution options that requires no account (TP2); the multilingual and accessibility infrastructure with a dedicated accessibility statement (TP4); an administrative create/read/update/delete area for the municipality; and, added in the most recent iteration as an additive contribution, a ``Get Involved'' page that groups the actions into four thematic initiatives via an accessible tab component, a section explaining how visitors can contribute (TP6), and a newsletter opt-in concept.

\textbf{Conceived (Future Work).} Persistent storage and moderation of freely submitted citizen initiatives; the actual sending of the newsletter; official UN~SDG iconography in place of coloured text badges; a dedicated school programme; an interactive map beyond simple location labels; and persistent gamification. These are deliberately \emph{not} built in the beta: introducing a production backend for personal data is a conscious design decision deferred to Future Work, where it must be handled with privacy-by-design (Diamantopoulou et al., 2019; Paguay-Chimarro et al., 2025).

% TODO: Bild einfügen — Screenshot der Startseite / Projektübersicht (Desktop + Mobile 375px)
\figplaceholder{Screenshot: ZOE landing page and action catalogue (desktop and 375\,px mobile)}

\subsection{Scenario~A --- Finding actions and contributing as a visitor (TP1, TP6)}

\emph{Persona:} a tourist staying near Acharavi for a week who wants to know how to help. The walkthrough begins on the action catalogue, switches to the ``Get Involved'' page, selects the \emph{Marine and Coast} initiative tab, and opens the ``How visitors can contribute'' section, which lists four concrete actions (join a beach cleanup, report litter, choose certified operators, respect wildlife). This demonstrates \textbf{DP1} (a central, structured, groupable place for all actions) and \textbf{DP6} (turning visitors from a burden into a resource). DP1 is justified by Simelio-Sol\`a et al. (2021), who show across 605 municipalities that unstructured municipal information prevents informed citizen decisions (Type~A); DP6 is justified by Vegas Macias et al. (2023) and Laksmi et al. (2026) (Type~A).

% TODO: Bild einfügen — Get-Involved-Seite mit Initiative-Tabs und Tourist:innen-Sektion
\figplaceholder{Screenshot: ``Get Involved'' page --- initiative tabs and visitor-contribution section}

\subsection{Scenario~B --- Submitting an idea without an account (TP2)}

\emph{Persona:} a resident of Kassiopi with an idea for a recurring cleanup. The walkthrough opens the participation page, chooses ``submit an idea'', completes the form \emph{without registering}, and receives confirmation. This demonstrates \textbf{DP2a} (lower the participation barrier by allowing contributions without a mandatory account). It is justified by Saldivar et al.\ (account/registration friction is a measurable barrier; Type~A) and by Arana-Catania et al. (2021) on information overload (Type~A). The walkthrough explicitly surfaces a limitation: in the prototype the submission is not persisted --- a deliberate scope decision, not an oversight. The associated point reward is presented with an evaluation caveat, because reward-based gamification alone can be short-lived (Thiel et al., 2016, Type~B); whether it raises or lowers perceived quality is a question for Phase~5.

% TODO: Bild einfügen — Beteiligungsformular (ohne Account) mit Bestätigung
\figplaceholder{Screenshot: participation form (no account) and confirmation state}

\subsection{Scenario~C --- Understanding SDG contribution and impact (TP3)}

\emph{Persona:} a teacher preparing a lesson who wants to know which SDGs the local actions serve and what they have achieved. The walkthrough moves from a project's SDG badges to the SDG dashboard and then to the transparency page with its impact indicators. This demonstrates \textbf{DP3a} (make SDG contribution and impact visible as a comprehensible dashboard) and \textbf{DP3b} (report impact in a SMART, comparable way and return the outcome to citizens). DP3a is justified by the experimental finding of Chokki et al.\ that well-designed dashboards increase engagement with open data (Type~A, preprint --- used with caution) and by Simonofski et al. (2022), who show that lay citizens expect simplified, visualised data (Type~A); the distinct status of \emph{outcome} transparency is established by Grimmelikhuijsen and Welch (2012, Type~A). The walkthrough makes explicit that all numeric values are fictional prototype data, clearly labelled as such.

% TODO: Bild einfügen — SDG-Dashboard und Transparenz-KPIs
\figplaceholder{Screenshot: SDG dashboard and transparency KPIs (impact values flagged as prototype data)}

\subsection{Scenario~D --- Reaching heterogeneous audiences (TP4)}

\emph{Persona:} a Greek-speaking elderly resident and, in contrast, a German-speaking visitor using a screen reader. The walkthrough switches the interface language between English, Greek, and German, navigates by keyboard only, and shows visible focus indicators and the accessibility statement. This demonstrates \textbf{DP4} (reach diverse groups through full multilingualism \emph{and} WCAG~2.1 AA accessibility in \emph{every} language version). It is justified by Csontos and Heckl (2021, 2025) on widespread WCAG non-conformance (Type~A), by Pontus and Rodr\'iguez V\'azquez (2021) on the accessibility gap between language versions (Type~A), and by Asghari et al. (2023) on the rarity of easy-language content (Type~A). The walkthrough is honest about a current limitation: several older pages still contain hard-coded English text and are therefore not yet fully multilingual; the new ``Get Involved'' page, by contrast, is fully internationalised across all three languages.

% TODO: Bild einfügen — Sprachumschaltung (EN/EL/DE) + Tastatur-Fokus-Indikator
\figplaceholder{Screenshot: language switch (EN/EL/DE) and keyboard focus indicator}

\section{Evaluation and Evaluability}

Demonstration is not evaluation; however, a central purpose of Phase~4 is to show that the artefact is \emph{evaluable} and to hand a concrete evaluation design to Phase~5. The Framework for Evaluation in Design Science (FEDS; Venable et al., 2016) structures this along two dimensions --- functional purpose (formative/summative) and paradigm (artificial/naturalistic) --- and recommends choosing an evaluation strategy explicitly. For ZOE, a \emph{Human Risk and Effectiveness} strategy is appropriate, because the principal uncertainty is whether real, heterogeneous citizens benefit.

\textbf{Artificial, ex-ante evaluation is already feasible.} Accessibility can be assessed both automatically (axe-core, already integrated into the test suite) and manually, following the method of Csontos and Heckl (2021, 2025); the criterion is WCAG~2.1 AA, the legal baseline being the EU Web Accessibility Directive 2016/2102 and EN~301~549. Front-end performance and quality can be audited with Lighthouse, exactly as Tsatsani et al. (2024) did for Greek municipal sites, yielding a comparable score. A multilingual completeness check (detecting missing translation keys across all three languages) operationalises the accessibility-across-languages concern of Pontus and Rodr\'iguez V\'azquez (2021). A structured heuristic walkthrough can additionally surface usability issues before any user is involved.

\textbf{Naturalistic, ex-post evaluation is planned.} Following FEDS, summative naturalistic episodes are designed but not yet executed: task-completion tests with representatives of each target group (``find action~X and its progress''; ``submit an idea without an account''; ``which SDGs does action~X address?''), the System Usability Scale, and comprehension and perceived-transparency measures. These are described per feature in the project's evaluation plan and handed over to Phase~5. The demonstration thus establishes the \emph{evaluability} that the next phase requires.

\section{Discussion, Limitations, and Future Work}

The demonstration shows that the design principles are realisable and behave as intended, supporting an affirmative answer to RQ1. It also makes the boundary between demonstration and evaluation explicit (RQ2): the scenarios show \emph{that} the artefact works, not yet \emph{how well} it works for real users.

Several limitations follow from deliberate design decisions. The free submission of citizen initiatives is not persisted, the newsletter is a concept without dispatch, SDG icons are coloured text badges rather than official iconography, and several legacy pages are not yet fully internationalised. These are conscious scope choices for a beta: a production backend that processes personal data is deferred to Future Work, where the transparency--privacy tension documented by Paguay-Chimarro et al. (2025) and the GDPR-compliance guidance of Diamantopoulou et al. (2019) must be respected. Other deliberate decisions --- the technology stack and the green visual identity (the stakeholder asked for a design ``greener than the existing municipal site'') --- carry no effect claim and are marked as such.

Future Work comprises the conceived features above, the sub-problems deferred as TP5/TP7/TP8, and persistent gamification. The project then transitions to Phase~6 (Communication), of which this report is itself a part.

\section{Conclusion}

This report has demonstrated a civic sustainability platform for North Corfu as the Phase~4 instantiation of a DSR project. Through four scenario walkthroughs, each tied to a design principle and to evidence classified as a measured effect or a framework recommendation, it has shown that the artefact addresses the visibility, participation, transparency, target-group, and tourism sub-problems, while honestly separating built from conceived features. By establishing that artificial evaluation is already feasible and that naturalistic evaluation is designed, it hands a concrete, FEDS-structured evaluation agenda to Phase~5.

\vspace{0.5em}
\small\emph{AI-use reflection.} Generative AI tools were used to support literature triage, drafting, and code scaffolding. All cited sources were verified against the original PDFs; uncertain bibliographic details are flagged in \texttt{docs/literature-index.md}. The author takes responsibility for the content.

\newpage
\section*{References}
\refitem{Arana-Catania, M., van Lier, F.-A., Procter, R., Tkachenko, N., He, Y., Zubiaga, A., \& Liakata, M. (2021). Citizen participation and machine learning for a better democracy. \emph{Digital Government: Research and Practice, 2}(3), Article~27. https://doi.org/10.1145/3452118}
\refitem{Asghari, H., Hewett, F., \& Z\"uger, T. (2023). On the prevalence of Leichte Sprache on the German web. In \emph{15th ACM Web Science Conference 2023 (WebSci~'23)}. ACM. https://doi.org/10.1145/3578503.3583599}
\refitem{Carayannis, E. G., Chatzichristofis, S., Kyrpianou, G., Psychalis, M., \& Sklias, G. (2026). The Greek smart city transition: Perceptions, barriers, and opportunities for sustainable urban growth. \emph{Journal of the Knowledge Economy, 17}, 6474--6501. https://doi.org/10.1007/s13132-025-02728-3}
\refitem{Chokki, A. P., Simonofski, A., Fr\'enay, B., \& Vanderose, B. (n.d.). Engaging citizens with open government data: The value of dashboards compared to individual visualizations. \emph{Digital Government: Research and Practice} (preprint).}
\refitem{Csontos, B., \& Heckl, I. (2021). Accessibility, usability, and security evaluation of Hungarian government websites. \emph{Universal Access in the Information Society, 20}, 139--156. https://doi.org/10.1007/s10209-020-00716-9}
\refitem{Csontos, B., \& Heckl, I. (2025). Five years of changes in the accessibility, usability, and security of Hungarian government websites. \emph{Universal Access in the Information Society, 24}, 2757--2781. https://doi.org/10.1007/s10209-025-01223-5}
\refitem{Diamantopoulou, V., Androutsopoulou, A., Gritzalis, S., \& Charalabidis, Y. (2019). Preserving digital privacy in e-participation environments: Towards GDPR compliance.}
\refitem{Gregor, S., \& Hevner, A. R. (2013). Positioning and presenting design science research for maximum impact. \emph{MIS Quarterly, 37}(2), 337--355. https://doi.org/10.25300/MISQ/2013/37.2.01}
\refitem{Gregor, S., \& Jones, D. (2007). The anatomy of a design theory. \emph{Journal of the Association for Information Systems, 8}(5), 312--335. https://doi.org/10.17705/1jais.00129}
\refitem{Grimmelikhuijsen, S. G., \& Welch, E. W. (2012). Developing and testing a theoretical framework for computer-mediated transparency of local governments. \emph{Public Administration Review.}}
\refitem{Hevner, A. R. (2007). A three cycle view of design science research. \emph{Scandinavian Journal of Information Systems, 19}(2), 87--92.}
\refitem{Laksmi, P. A. S., Suriani, N. N., Arjawa, I. G. W., \& Saputra, K. A. K. (2026). Sustainable tourism development in the digital era: Assessing the impact of competitive strategies and community participation. \emph{European Journal of Sustainable Development Research, 10}(2), em0373. https://doi.org/10.29333/ejosdr/17841}
\refitem{Peffers, K., Tuunanen, T., Rothenberger, M. A., \& Chatterjee, S. (2007). A design science research methodology for information systems research. \emph{Journal of Management Information Systems, 24}(3), 45--77. https://doi.org/10.2753/MIS0742-1222240302}
\refitem{Pontus, V., \& Rodr\'iguez V\'azquez, S. (2021). Language-related criteria for evaluating the accessibility of localised multilingual websites. In \emph{Proceedings of the 3rd Swiss Conference on Barrier-free Communication (BfC~2020)} (pp.~162--169). ZHAW.}
\refitem{Royo, S., Pina, V., \& Garcia-Rayado, J. (2020). Decide Madrid: A critical analysis of an award-winning e-participation initiative. \emph{Sustainability, 12}(4), 1674. https://doi.org/10.3390/su12041674}
\refitem{Royo, S., Bell\`o, B., Torres, L., \& Downe, J. (2024). The success of e-participation: Learning lessons from Decide Madrid and We Asked, You Said, We Did in Scotland. \emph{Policy \& Internet, 16}(1), 65--82. https://doi.org/10.1002/poi3.363}
\refitem{Saldivar, J., Daniel, F., Cernuzzi, L., \& Casati, F. (n.d.). Online idea management for civic engagement: A study on the benefits of integration with social networking. \emph{Computer Supported Cooperative Work (CSCW).}}
\refitem{Simelio-Sol\`a, N., Ferr\'e-Pavia, C., \& Herrero-Guti\'errez, F.-J. (2021). Transparent information and access to citizen participation on municipal websites. \emph{Profesional de la Informaci\'on, 30}(2), e300211. https://doi.org/10.3145/epi.2021.mar.11}
\refitem{Simonofski, A., Zuiderwijk, A., Clarinval, A., \& Hammedi, W. (2022). Tailoring open government data portals for lay citizens: A gamification theory approach. \emph{International Journal of Information Management, 65}, 102511. https://doi.org/10.1016/j.ijinfomgt.2022.102511}
\refitem{Thiel, S.-K., Reisinger, M., R\"oderer, K., \& Fr\"ohlich, P. (2016). Playing (with) democracy: A review of gamified participation approaches. \emph{JeDEM, 8}(3), 32--60.}
\refitem{Tsatsani, O., Patergiannaki, Z., \& Pollalis, Y. (2024). Unraveling the dynamics of e-government digitization, penetration and user experience: A case study of Greek municipalities. \emph{Journal of Strategic Innovation and Sustainability, 19}(3).}
\refitem{Vegas Macias, J., Lamers, M., \& Toonen, H. (2023). Connecting fishing and tourism practices using digital tools: A case study of Marsaxlokk, Malta. \emph{Maritime Studies, 22}, 24. https://doi.org/10.1007/s40152-023-00305-5}
\refitem{Venable, J., Pries-Heje, J., \& Baskerville, R. (2016). FEDS: A framework for evaluation in design science research. \emph{European Journal of Information Systems, 25}(1), 77--89. https://doi.org/10.1057/ejis.2014.36}

\end{document}
